% New research note, 16 October 2026. Uses the established pseudo-unitary
% reduction and Ising projection in frontier_1660.tex. Not a categorification.
\subsection{The complete twist-one support at multiplicity eight}
\label{subsec:1660-twist-one-support}

The equality case of the dimension-block bound gives a useful sharpening
of the center reduction. It also gives a shorter proof of tensor-quotient
rigidity for this ring. It does not decide categorifiability.

\begin{proposition}\label{prop:1660-twist-one-support}
Let $\mathcal C$ be a pseudo-unitary categorification of $\RR{8}{1660}$,
with its canonical positive spherical structure. The twist-one simple
objects of $Z(\mathcal C)$ are exactly the six simple summands of
$I_{\mathcal C}(1)$.
\end{proposition}
\begin{proof}
Put $d=3+2\sqrt2+2\sqrt3+\sqrt6$ and
$D=4(1+d^2)=144+96\sqrt2+80\sqrt3+56\sqrt6$. Then
\[
 \frac1D=\frac{6-2\sqrt3-\sqrt6}{48}.
\]
The dimension and global-dimension fields are both
$K=\mathbb Q(\sqrt2,\sqrt3)$, so the dimensional grading has order one.
For the characters $\chi_3,\chi_6$ of $\operatorname{Gal}(K/\mathbb Q)$
corresponding to $\sqrt3,\sqrt6$, the Fourier coefficients of $1/D$ are
$-\sqrt3/6$ and $-\sqrt6/12$. Their squared absolute values are,
respectively, $1/12$ and $1/24$, equal to the reciprocals of their
conductors. Thus both attain equality in the dimension-block bound
of Schopieray, \emph{The dimension block of a spherical fusion category},
Theorem~5.1, arXiv:2609.35287.

Here is the equality consequence we use. Let $L_\gamma$, for
$\gamma\in\operatorname{Gal}(K/\mathbb Q)$, be the four summands of
$I(1)$ indexed by the conjugates of the dimension character, and set
$w_\chi=\sum_\gamma\chi(\gamma)e_{L_\gamma}$. Write $s$ for the normalized modular matrix, and let $P$ be the orthogonal
projection onto the twist-one eigenspace of its modular representation. The proof of the cited theorem gives
\[
 \|Psw_\chi\|^2\leq\frac{4}{\mathfrak f(\chi)}.
\]
The four dimension-block coordinates already contribute exactly
$4|\widehat{(1/D)}(\chi)|^2=4/\mathfrak f(\chi)$.
Consequently every other twist-one simple $Z$ has zero $\chi$-Fourier
coefficient in $\dim(Z)/D$. Applying this to $\chi_3,\chi_6$ gives
\[
 \theta_Z=1,\quad Z\notin\{L_\gamma\}
 \quad\Longrightarrow\quad \frac{\dim Z}{D}\in\mathbb Q(\sqrt2).
 \tag{*}
\]
This uses unitarity of the modular representation, not a unitary
structure on $\mathcal C$.

Suppose that $Z$ is twist-one and not a summand of $I(1)$.
Adjunction gives $[FZ:1]=0$. Its $g$-multiplicity is also zero:
the established central fermion $h$ gives $I(g)=hI(1)$; the four even
low constituents have $h$-translates of twist $-1$, while the two odd
low constituents are permuted among themselves by $h$.
Therefore write
\[
 FZ=bX+c\,gX+z\sigma+aT,\qquad T=\sigma X.
\]
Since $\dim Z=z\sqrt2+(b+c+a\sqrt2)d$, direct expansion of $(*)$ gives
\[
 b+c=4z,\qquad a=3z.
\]
If $z=0$, this is the zero restriction vector, impossible for a
nonzero object. If $z>0$, the Ising projection established in
Proposition~\ref{prop:1660-uniform-indicators} forces the twist of $Z$
to have primitive order $16$, again a contradiction.
Finally, all six summands of $I(1)$ have twist one, and each has
multiplicity one because $K_0(\mathcal C)$ is commutative.
\end{proof}

\begin{corollary}\label{cor:1660-twist-one-quotient}
Every tensor functor from a complex categorification of
$\RR{8}{1660}$ to a complex fusion category is fully faithful.
\end{corollary}
\begin{proof}
Replace the target by the fusion subcategory generated by the image,
so that the functor is dominant. Apply the established Galois reduction,
and use the established pseudo-unitary descent lemma to give source
and target their canonical spherical structures; the functor preserves
these structures. Its induced central algebra $A$ is a connected
\'etale algebra of twist one in $Z(\mathcal C)$, and its underlying
algebra is connected: if $R$ is the right adjoint, then
\[
 \dim\operatorname{Hom}_{\mathcal C}(1,R(1))
 =\dim\operatorname{Hom}_{\mathcal D}(1,1)=1.
\]
By the proposition every constituent of $A$ lies in $I(1)$.
Every such simple restricts with unit multiplicity one. Since the
unit already occurs in $A$, underlying connectedness forces $A=1$.
The projection formula gives $RF(V)=V\otimes R(1)=V$, so $F$ is
fully faithful, and hence an equivalence onto its image. Transporting
back under the Galois conjugation proves the assertion.
\end{proof}
